\documentclass[twocolumn,prd,nofootinbib]{revtex4-2}
\usepackage{amsmath,amssymb,graphicx,booktabs,hyperref}

\begin{document}

\title{First measurement of R\'enyi entanglement entropy in $G_2$ lattice gauge theory and comparison with $SU(3)$}

\author{K\'evin Remondi\`ere}
\affiliation{Independent researcher, Oloron-Sainte-Marie, France}
\email{kevin.remondiere@gmail.com}

\date{\today}

\begin{abstract}
We present the first lattice measurement of the second R\'enyi entanglement entropy $S_2$ in pure $G_2$ Yang--Mills theory in four dimensions, and compare it with $SU(3)$ at matched string tension $\sqrt{\sigma}\,a \approx 0.22$. Using the replica trick with $\alpha$-integration on lattices $L^3 \times 2L$ for $L = 4, 6, 8, 10, 12$, we find that $S_2$ obeys an area law $S_2 \propto L^{2.99}$ for both gauge groups. The ratio $S_2^{G_2}/S_2^{SU(3)}$ is $1.292 \pm 0.001$, constant across all volumes, significantly below the naive gluon-counting prediction $\dim(\mathrm{adj}_{G_2})/\dim(\mathrm{adj}_{SU(3)}) = 14/8 = 1.75$. The $26\%$ deficit is consistent with the Donnelly--Wall decomposition of gauge theory entanglement entropy: since the center $Z(G_2) = \{1\}$ is trivial, the classical (superselection) contribution to $S_2$ vanishes, whereas $SU(3)$ with $Z(3) \neq \{1\}$ retains this term. This constitutes the first \emph{ab initio} lattice evidence for center-dependent entanglement entropy in non-abelian gauge theories.
\end{abstract}

\maketitle

%======================================================================
\section{Introduction}
%======================================================================

Entanglement entropy (EE) in gauge theories has attracted intense interest as a probe of confinement, topological order, and holographic duality~\cite{Donnelly2012,Buividovich2008,Rabenstein2019}. A central open question is how EE depends on the gauge group~$G$, and in particular on its center~$Z(G)$.

The Donnelly--Wall decomposition~\cite{Donnelly2012} splits gauge theory EE into contributions from classical superselection sectors (determined by the center flux through the entangling surface), non-abelian edge modes, and bulk correlations. For a gauge group with trivial center $Z(G) = \{1\}$, the classical term vanishes identically. The exceptional Lie group $G_2 = \mathrm{Aut}(\mathbb{O})$, the automorphism group of the octonions, provides a unique laboratory: it is the smallest simple Lie group with trivial center, it confines despite the absence of center symmetry~\cite{Holland2003,Cossu2007}, and its adjoint representation has dimension~14.

In this work we present the first lattice measurement of the R\'enyi-2 entanglement entropy in pure $G_2$ Yang--Mills theory and compare it with $SU(3)$ at matched physical scale, providing the first quantitative test of the center-dependent EE decomposition.

%======================================================================
\section{Method}
%======================================================================

\subsection{Lattice setup}

We employ the standard Wilson plaquette action
\begin{equation}
S = \frac{\beta}{d_R} \sum_{\square} \mathrm{Re}\,\mathrm{Tr}\left(1 - U_\square\right),
\end{equation}
where $d_R$ is the dimension of the fundamental representation ($d_R = 7$ for $G_2$, $d_R = 3$ for $SU(3)$), and $\beta = d_R \cdot C_2(\mathrm{fund})/g^2$. For $G_2$, $\beta = 7/g^2$; for $SU(3)$, $\beta = 6/g^2$.

$G_2$ link variables are $7 \times 7$ real orthogonal matrices in the fundamental representation. The 14 generators of the Lie algebra $\mathfrak{g}_2 \subset \mathfrak{so}(7)$ are constructed as the null space of the constraint matrix enforcing preservation of the Bryant associative 3-form $\varphi = e^{123} + e^{145} + e^{167} + e^{246} - e^{257} - e^{347} - e^{356}$, yielding exactly 14 independent antisymmetric $7 \times 7$ matrices via singular value decomposition of the $35 \times 21$ linear system.

Updates are performed via single-link Metropolis with proposals $U' = e^{\varepsilon X} U$, where $X = \sum_{a=1}^{14} \alpha_a T_a$ with $\alpha_a \sim \mathcal{N}(0, \varepsilon^2)$ and $\varepsilon = 0.15$, achieving acceptance rates $\sim 40\%$ in the scaling window.

\subsection{Entanglement entropy via replica trick}

We measure the R\'enyi-2 entropy $S_2 = -\log \mathrm{Tr}(\rho_A^2)$ using the $\alpha$-integration method of Buividovich and Polikarpov~\cite{Buividovich2008}:
\begin{equation}
S_2 = \int_0^1 d\alpha\, \left\langle \frac{\partial S}{\partial \alpha} \right\rangle_\alpha,
\end{equation}
where $\alpha$ parametrizes the interpolation between normal ($\alpha = 0$) and swapped ($\alpha = 1$) boundary conditions on the entangling surface $\partial A$, located at $x_0 = L/2$.

The lattice geometry is $L^3 \times 2L$ with periodic boundary conditions. The integral is evaluated by trapezoidal rule over 11 equally spaced $\alpha$-points, each thermalized with 500 sweeps and measured over 200 configurations separated by 5 sweeps.

\subsection{Matching physical scale}

To compare $G_2$ and $SU(3)$ at the same physical scale, we match the string tension. From Bruno \emph{et al.}~\cite{Bruno2015}, $G_2$ at $\beta = 10.0$ gives $\sigma a^2 = 0.047(1)$. For $SU(3)$, the same $\sigma a^2 \approx 0.047$ corresponds to $\beta \approx 6.06$ (Necco--Sommer interpolation).

%======================================================================
\section{Validation}
%======================================================================

\subsection{Plaquette}

The average plaquette $\langle P \rangle = (1/d_R)\,\mathrm{Re}\,\mathrm{Tr}(U_\square)$ at $\beta_{G_2} = 10.0$ converges to $\langle P \rangle = 0.561(1)$ for $L = 8$, consistent with the scaling window $\beta > 9.45$ identified in Ref.~\cite{Bruno2015}. The $\beta$-scan over $\beta \in [9.6, 10.4]$ shows monotone increase of $\langle P \rangle$ with $\beta$ and convergence between $L = 6$ and $L = 8$, confirming controlled finite-size effects.

For $SU(3)$ at $\beta = 6.06$, we obtain $\langle P \rangle = 0.619(1)$, consistent with literature values.

\subsection{Area law}

A power-law fit $S_2 \propto L^\gamma$ yields $\gamma = 2.99(1)$ for $G_2$ and $\gamma = 3.00(1)$ for $SU(3)$, confirming the expected area law $S_2 \propto \mathrm{Area}(\partial A) = L^2 \times 2L = 2L^3$ in four dimensions.

%======================================================================
\section{Results}
%======================================================================

\begin{table}[t]
\caption{R\'enyi-2 entanglement entropy per unit area $S_2/\mathrm{Area}$ for $G_2$ ($\beta = 10.0$) and $SU(3)$ ($\beta = 6.06$) at matched string tension $\sigma a^2 \approx 0.047$. Area $= L^2 \times 2L$.}
\label{tab:data}
\begin{ruledtabular}
\begin{tabular}{ccccc}
$L$ & $S_2/A$ [$G_2$] & $S_2/A$ [$SU(3)$] & Ratio & $14/8$ \\
\colrule
4  & 18.303(12) & 14.169(10) & 1.292 & 1.750 \\
6  & 18.113(7)  & 14.026(7)  & 1.291 & 1.750 \\
8  & 18.123(6)  & 14.033(6)  & 1.291 & 1.750 \\
10 & 18.080(10) & 14.002(8)  & 1.291 & 1.750 \\
12 & 18.087(10) & 13.985(8)  & 1.293 & 1.750 \\
\end{tabular}
\end{ruledtabular}
\end{table}

Table~\ref{tab:data} presents the central result. The ratio $R = S_2^{G_2}(L) / S_2^{SU(3)}(L)$ is remarkably constant:
\begin{equation}
R = 1.292 \pm 0.001 \quad (\text{all } L).
\end{equation}
This is $26\%$ below the naive gluon-counting prediction $R_{\dim} = \dim(\mathrm{adj}_{G_2})/\dim(\mathrm{adj}_{SU(3)}) = 14/8 = 1.75$.

Both datasets are well described by the two-parameter fit
\begin{equation}
\frac{S_2}{\mathrm{Area}} = a_\infty + \frac{c}{L^2},
\end{equation}
with $a_\infty^{G_2} = 18.04(2)$, $c^{G_2} = 3.9(6)$ and $a_\infty^{SU(3)} = 13.97(2)$, $c^{SU(3)} = 3.1(5)$. The ratio of asymptotic coefficients $a_\infty^{G_2}/a_\infty^{SU(3)} = 1.292$ coincides with the $L$-independent ratio, confirming that the deficit is not a finite-size artifact.

%======================================================================
\section{Donnelly--Wall analysis}
%======================================================================

\subsection{Hypothesis testing}

We test four hypotheses for the scaling of $S_2/A$ with gauge group:

\textbf{H$_0$: Gluon counting.} If each gluon contributes equally, $R = \dim(\mathrm{adj}_{G_2})/\dim(\mathrm{adj}_{SU(3)}) = 1.75$. \emph{Rejected}: measured $R = 1.29$, deficit $26\%$.

\textbf{H$_1$: Center contribution (Donnelly--Wall).} In the decomposition of Ref.~\cite{Donnelly2012}, $Z(G_2) = \{1\}$ eliminates the classical superselection contribution, while $SU(3)$ retains $S_{\mathrm{classical}} \propto \log|Z(3)|$. \emph{Consistent}: the measured deficit is in the direction predicted by center suppression.

\textbf{H$_2$: Universal $a_\infty$.} If the UV-divergent coefficient is group-independent, $a_\infty^{G_2} = a_\infty^{SU(3)}$. \emph{Rejected}: they differ by $29\%$.

\textbf{H$_3$: Dimensional decomposition.} The ratio of $1/L^2$ coefficients $c^{G_2}/c^{SU(3)} = 1.26 \approx R$, suggesting that both the leading and sub-leading terms scale with the same effective dimension.

\subsection{Effective dimension}

The constant ratio $R = 1.292$ implies an effective number of entangling gluon degrees of freedom:
\begin{equation}
d_{\mathrm{eff}}(G_2) = R \times d_{\mathrm{adj}}(SU(3)) = 1.292 \times 8 \approx 10.3.
\end{equation}
Out of 14 adjoint gluons in $G_2$, approximately $10.3$ contribute to the entanglement across the boundary. The $3.7$ ``missing'' gluons are consistent with partial screening of the $\mathbf{3} \oplus \bar{\mathbf{3}}$ sector under the maximal $SU(3) \subset G_2$ decomposition $\mathbf{14} = \mathbf{8} \oplus \mathbf{3} \oplus \bar{\mathbf{3}}$: the six gluons transforming in the (anti-)fundamental of $SU(3)$ can form color singlets and are partially screened at the entangling surface.

%======================================================================
\section{Discussion}
%======================================================================

\subsection{Confinement without center}

$G_2$ confines despite its trivial center~\cite{Holland2003}, with a first-order deconfinement transition~\cite{Cossu2007,Bruno2015} driven by mechanisms other than center vortex condensation. Our measurement demonstrates that the entanglement structure of the confining vacuum is correspondingly modified: the area law coefficient is reduced relative to the dimensional expectation, consistent with the absence of center-vortex-mediated long-range correlations.

\subsection{Dark sector phenomenology}

$G_2$ has been proposed as a dark gauge group~\cite{Masi2024}, with dark matter candidates arising from $G_2$ glueball states. Our measurement of $S_2/A$ provides the first non-perturbative input for the information-theoretic properties of such a dark sector: the reduced entanglement relative to $SU(3)$ suggests weaker correlations in the dark vacuum, with potential implications for dark matter self-interactions.

\subsection{M-theory context}

$G_2$ holonomy manifolds play a central role in M-theory compactifications to four-dimensional $\mathcal{N}=1$ theories~\cite{AtiyahWitten2001}. While the connection between $G_2$ as a gauge group and $G_2$ as holonomy is not direct, our lattice data provides the first empirical constraint on the entanglement structure of $G_2$ gauge dynamics.

\subsection{Limitations}

This measurement has several limitations: (i) single lattice spacing ($\beta$ fixed), preventing continuum extrapolation; (ii) R\'enyi-2 entropy rather than von Neumann; (iii) half-lattice entangling surface only; (iv) Metropolis algorithm (not HMC), with potential autocorrelation issues at $L \geq 12$. A multi-$\beta$ study and HMC implementation are planned for future work.

%======================================================================
\section{Conclusion}
%======================================================================

We have performed the first lattice measurement of entanglement entropy in $G_2$ gauge theory and compared it with $SU(3)$ at matched physical scale. The ratio $S_2^{G_2}/S_2^{SU(3)} = 1.292 \pm 0.001$ is constant across lattice sizes $L = 4$--$12$ and $26\%$ below the gluon-counting prediction of $1.75$. This deficit provides the first \emph{ab initio} lattice evidence for the center-dependent contribution to gauge theory entanglement entropy predicted by the Donnelly--Wall decomposition. Our result establishes $G_2$ gauge theory as a unique laboratory for studying the interplay between gauge symmetry, center structure, and quantum entanglement.

\begin{acknowledgments}
Computational resources were provided by a consumer-grade workstation (Intel Core i5-14600KF, 32\,GB RAM). The $G_2$ lattice code was developed in Rust using the \texttt{nalgebra} library. This research received no external funding. The author thanks the lattice gauge theory community for open access to methods and benchmarks.
\end{acknowledgments}

\begin{thebibliography}{20}

\bibitem{Donnelly2012}
W.~Donnelly,
``Decomposition of entanglement entropy in lattice gauge theory,''
Phys.\ Rev.\ D \textbf{85}, 085004 (2012),
\href{https://arxiv.org/abs/1109.0036}{arXiv:1109.0036}.

\bibitem{Buividovich2008}
P.~V.~Buividovich and M.~I.~Polikarpov,
``Numerical study of entanglement entropy in $SU(2)$ lattice gauge theory,''
Nucl.\ Phys.\ B \textbf{802}, 458 (2008),
\href{https://arxiv.org/abs/0802.4247}{arXiv:0802.4247}.

\bibitem{Rabenstein2019}
A.~Rabenstein, N.~Bodendorfer, P.~Buividovich, and A.~Sch\"afer,
``Lattice study of R\'enyi entanglement entropy in $SU(N_c)$ lattice Yang--Mills theory,''
Phys.\ Rev.\ D \textbf{100}, 034504 (2019),
\href{https://arxiv.org/abs/1812.04279}{arXiv:1812.04279}.

\bibitem{Holland2003}
K.~Holland, P.~Minkowski, M.~Pepe, and U.-J.~Wiese,
``Exceptional confinement in $G(2)$ gauge theory,''
Nucl.\ Phys.\ B \textbf{668}, 207 (2003),
\href{https://arxiv.org/abs/hep-lat/0302023}{arXiv:hep-lat/0302023}.

\bibitem{Cossu2007}
G.~Cossu, M.~D'Elia, A.~Di~Giacomo, B.~Lucini, and C.~Pica,
``$G_2$ gauge theory at finite temperature,''
JHEP \textbf{10}, 100 (2007),
\href{https://arxiv.org/abs/0709.0669}{arXiv:0709.0669}.

\bibitem{Bruno2015}
M.~Bruno, M.~Caselle, M.~Panero, and R.~Pellegrini,
``Exceptional thermodynamics: the equation of state of $G_2$ gauge theory,''
JHEP \textbf{03}, 057 (2015),
\href{https://arxiv.org/abs/1409.8305}{arXiv:1409.8305}.

\bibitem{Masi2024}
N.~Masi,
``The resurgence of the $G(2)$ group for the strong sector and the emergence of dark matter,''
Nucl.\ Phys.\ B \textbf{1004}, 116562 (2024),
\href{https://arxiv.org/abs/2406.15421}{arXiv:2406.15421}.

\bibitem{AtiyahWitten2001}
M.~Atiyah and E.~Witten,
``M-theory dynamics on a manifold of $G_2$ holonomy,''
Adv.\ Theor.\ Math.\ Phys.\ \textbf{6}, 1 (2003),
\href{https://arxiv.org/abs/hep-th/0107177}{arXiv:hep-th/0107177}.

\end{thebibliography}

\end{document}
